\documentclass[10pt,a4paper]{amsart}
\usepackage{amsmath,amssymb,amsthm,graphicx}
\usepackage{verbatim}
\newenvironment{reference}{\comment}{\endcomment}
%\newenvironment{reference}{}{}
\newenvironment{slogan}{\comment}{\endcomment}
\newenvironment{history}{\comment}{\endcomment}

% For commutative diagrams we use Xy-pic
\usepackage[all]{xy}

% We use 2cell for 2-commutative diagrams.
\xyoption{2cell}
% Theorem environments.
%
\theoremstyle{plain}
\newtheorem{theorem}[subsection]{Theorem}
\newtheorem{proposition}[subsection]{Proposition}
\newtheorem{lemma}[subsection]{Lemma}

\theoremstyle{definition}
\newtheorem{definition}[subsection]{Definition}
\newtheorem{example}[subsection]{Example}
\newtheorem{exercise}[subsection]{Exercise}
\newtheorem{situation}[subsection]{Situation}

\theoremstyle{remark}
\newtheorem{remark}[subsection]{Remark}
\newtheorem{remarks}[subsection]{Remarks}

\numberwithin{equation}{subsection}

% Macros
%
\def\lim{\mathop{\mathrm{lim}}\nolimits}
\def\colim{\mathop{\mathrm{colim}}\nolimits}
\def\Spec{\mathop{\mathrm{Spec}}}
\def\Hom{\mathop{\mathrm{Hom}}\nolimits}
\def\Ext{\mathop{\mathrm{Ext}}\nolimits}
\def\SheafHom{\mathop{\mathcal{H}\!\mathit{om}}\nolimits}
\def\SheafExt{\mathop{\mathcal{E}\!\mathit{xt}}\nolimits}
\def\Sch{\mathit{Sch}}
\def\Mor{\mathop{\mathrm{Mor}}\nolimits}
\def\Ob{\mathop{\mathrm{Ob}}\nolimits}
\def\Sh{\mathop{\mathit{Sh}}\nolimits}
\def\NL{\mathop{N\!L}\nolimits}
\def\CH{\mathop{\mathrm{CH}}\nolimits}
\def\proetale{{pro\text{-}\acute{e}tale}}
\def\etale{{\acute{e}tale}}
\def\QCoh{\mathit{QCoh}}
\def\Ker{\mathop{\mathrm{Ker}}}
\def\Im{\mathop{\mathrm{Im}}}
\def\Coker{\mathop{\mathrm{Coker}}}
\def\Coim{\mathop{\mathrm{Coim}}}

% Boxtimes
%
\DeclareMathSymbol{\boxtimes}{\mathbin}{AMSa}{"02}

%
% Macros for moduli stacks/spaces
%
\def\QCohstack{\mathcal{QC}\!\mathit{oh}}
\def\Cohstack{\mathcal{C}\!\mathit{oh}}
\def\Spacesstack{\mathcal{S}\!\mathit{paces}}
\def\Quotfunctor{\mathrm{Quot}}
\def\Hilbfunctor{\mathrm{Hilb}}
\def\Curvesstack{\mathcal{C}\!\mathit{urves}}
\def\Polarizedstack{\mathcal{P}\!\mathit{olarized}}
\def\Complexesstack{\mathcal{C}\!\mathit{omplexes}}
% \Pic is the operator that assigns to X its picard group, usage \Pic(X)
% \Picardstack_{X/B} denotes the Picard stack of X over B
% \Picardfunctor_{X/B} denotes the Picard functor of X over B
\def\Pic{\mathop{\mathrm{Pic}}\nolimits}
\def\Picardstack{\mathcal{P}\!\mathit{ic}}
\def\Picardfunctor{\mathrm{Pic}}
\def\Deformationcategory{\mathcal{D}\!\mathit{ef}}
\usepackage{cleveref}
\crefname{originalreference}{original reference}{original references}
\Crefname{originalreference}{Original reference}{Original references}
\makeatletter
\newlabel{stacks-properties-lemma-topology-points}{{04XL}{source}{}{sourcepointer-1}{}}
\newlabel{stacks-properties-lemma-topology-points@cref}{{[originalreference][0][]04XL}{[1][0][]source}{}{}{}}
\makeatother
\title{Images of irreducible components under smooth morphisms are irreducible components; surjective case exhausts}
\author{The Stacks Project Authors}
\date{stacks-project:stacks-geometry.tex:local-snapshot-20261012; GNU FDL 1.2 or later}
\begin{document}
\maketitle
\noindent\textit{Editorial note.} All delivered mathematical text is the author's own, copied verbatim from the pinned Stacks Project chapter file stacks-geometry.tex; no statement or proof step was rewritten, reordered or removed.\par\medskip
\noindent\textit{Editorial note.} Span recomputation by label pairing (LANE\_BRIEF 10.2): in this chapter the author's proof environment follows the statement environment as a standalone block; the spec covers the complete environment pair (statement span, proof span) as recorded in the pool candidate.\par\medskip
\noindent\textit{Editorial note.} Occupation scan: combined claim\_key + statement\_anchor + generic-line scan over all registration files (active and withdrawn) found 0 hits for stacks-geometry.tex (ge-occupied.json, 2026-10-14); 207 route=geometry / chapter-list field hits in focused-seeds are alias noise (route name and alphabetical chapter enumerations), excluded per the ladder alias-blind-spot caveat.\par\medskip
\ifdefined\hypertarget\hypertarget{sourcepointer-1}{}\fi
\noindent\textit{Source reference.} Original reference 04XL: sources/stacks-chow/tags.txt, source lines 6399--6399. Cross-chapter standard prerequisite cited by the author inside the counted span; it is used as a pointer and its own proof is not part of this item.\par\medskip
\begin{lemma}
\label{lemma-map-of-components}
Let $f : U \to \mathcal{X}$ be a smooth morphism from a scheme
to a locally Noetherian algebraic stack. The closure of the image of any
irreducible component of $|U|$ is an irreducible component of $|\mathcal{X}|$.
If $U \to \mathcal{X}$ is surjective, then all irreducible components of
$|\mathcal{X}|$ are obtained in this way.
\end{lemma}

\begin{proof}
The map $|U| \to |\mathcal{X}|$ is continuous and open by
Properties of Stacks, Lemma \ref{stacks-properties-lemma-topology-points}.
Let $T \subset |U|$ be an irreducible component. Since $U$ is locally
Noetherian, we can find a nonempty affine open $W \subset U$ contained in $T$.
Then $f(T) \subset |\mathcal{X}|$ is irreducible and contains the
nonempty open subset $f(W)$. Thus the closure of $f(T)$ is irreducible and
contains a nonempty open. It follows that this closure is an irreducible
component.

\medskip\noindent
Assume $U \to \mathcal{X}$ is surjective and let $Z \subset |\mathcal{X}|$
be an irreducible component. Choose a Noetherian open subset $V$
of $|\mathcal{X}|$ meeting $Z$. After removing the other irreducible
components from $V$ we may assume that $V \subset Z$.
Take an irreducible component of the nonempty
open $f^{-1}(V) \subset |U|$ and let $T \subset |U|$ be its closure.
This is an irreducible component of $|U|$ and the closure of $f(T)$
must agree with $Z$ by our choice of $T$.
\end{proof}

\end{document}
